%% =====================================================================
%%  Unit 8 — Home sweet home  (Décrire son logement)
%%  Pistes : 8.01–8.04 dialogues · 8.05–8.06 vocabulaire · 8.07–8.08 grammaire
%%           8.09 prononciation · 8.10–8.13 exercices (textes en annexe)
%% =====================================================================

\newcommand{\otIllusUnVIII}{\otIllusDuo{kitchen/104/86/15/0.10,sofa/190/92/20/0.09,bedroom/122/14/13/0.08,
  window/150/8/11/0.08,door/196/22/12/0.07,home/100/50/12/0.07}{emma}{Is there a bathroom?}{lea}{Yes, there is!}}

\unite[Home sweet\\home]{8}{Home sweet home}{Décrire son logement}{otIllusUnVIII}

\situation{Épisode 8}{Dimanche après-midi. Pour la première fois, Léa reçoit Emma dans son petit
appartement de Brooklyn : visite guidée ! Le lendemain matin, catastrophe : Léa est en retard et
son téléphone a disparu. Heureusement, Sam passe par là.}

\begin{objectifs}
\item décrire votre logement, pièce par pièce ;
\item dire ce qu'il y a (ou ce qu'il n'y a pas) avec \en{there is} et \en{there are} ;
\item situer un objet : \en{in}, \en{on}, \en{under}, \en{next to}, \en{behind}… ;
\item demander où se trouve quelque chose : \en{Where's…?} \en{Where are…?} ;
\item prononcer le son \phon{θ} de \en{three} et \en{bathroom}.
\end{objectifs}

\auprogramme
  {Décrire un logement, situer des objets}
  {\en{there is}, \en{there are}\par\vspace{0.6mm}Les prépositions de lieu}
  {Les pièces, les meubles : 27 expressions}
  {Le son \phon{θ} de \en{three}}
  {\en{Apartment} ou \en{flat} ? Se loger à New York}

\semaine{casque/otNavy/Dialogues/Écouter et lire les deux scènes,
  bulle/otGreenDark/Vocabulaire/Apprendre les 27 expressions,
  blocs/otViolet/Grammaire/there is et prépositions,
  sons/otOrangeDark/Sons et culture/Le son de three et les logements,
  crayon/otNavy/Exercices/12 exercices corrigés,
  coche/otGreenDark/Bilan/L'essentiel et le test express}

\clearpage
%% ---------------------------------------------------------------------
\rubrique{ecouter}{Chez Léa}

\begin{dialogue}{1}{Welcome to my apartment!}{Dimanche, 15 heures, à Brooklyn. Emma arrive chez Léa
  pour la première fois.}{8.01}{8.02}
\dit{lea}{Welcome to my apartment!}{Bienvenue dans mon appartement !}
\dit{emma}{Thank you! Oh, it's lovely!}{Merci ! Oh, il est charmant !}
\dit{lea}{\hl{There's a} small living room, \hl{there's a} kitchen and \hl{there's a} bedroom.\note{1}}{Il y a un petit salon, une cuisine et une chambre.}
\dit{emma}{\hl{Is there a} bathroom?}{Il y a une salle de bains ?}
\dit{lea}{\hl{Yes, there is}. It's \hl{next to} the bedroom.\note{2}}{Oui. Elle est à côté de la chambre.}
\dit{emma}{Wow, \hl{there are} two big windows!}{Waouh, il y a deux grandes fenêtres !}
\dit{lea}{Yes, and look: there are plants \hl{on} the fire escape!\note{3}}{Oui, et regarde : il y a des plantes sur l'escalier de secours !}
\dit{emma}{But \hl{there isn't a} washing machine?}{Mais il n'y a pas de machine à laver ?}
\dit{lea}{\hl{No, there isn't}. There's a laundromat \hl{on the corner}.}{Non. Il y a une laverie au coin de la rue.}
\dit{emma}{That's very New York!}{C'est très new-yorkais !}
\end{dialogue}

\begin{comprendre}
\pointnote{1} \en{\textbf{There's a kitchen.}} « Il y a une cuisine. » \en{There's} = \en{There
  is}, devant un nom singulier. Devant un pluriel : \en{There are two windows}.
\pointnote{2} \en{\textbf{next to the bedroom}} « à côté de la chambre ». \en{next to} est une
  préposition de lieu, comme \en{in} (dans), \en{on} (sur) ou \en{under} (sous).
\pointnote{3} \en{\textbf{the fire escape}} Les escaliers métalliques fixés aux façades des
  immeubles new-yorkais, en cas d'incendie. Beaucoup d'habitants y installent des plantes !
\end{comprendre}

\begin{dialogue}{2}{Where's my phone?}{Lundi, 8 h 30. Léa doit partir au studio, mais elle ne
  trouve plus son téléphone. Sam frappe à la porte.}{8.03}{8.04}
\dit{lea}{Oh no! \hl{Where's} my phone?\note{4}}{Oh non ! Où est mon téléphone ?}
\dit{sam}{Is it \hl{in} your bag?}{Il est dans ton sac ?}
\dit{lea}{No, it isn't.}{Non.}
\dit{sam}{Is it \hl{on} the table?}{Il est sur la table ?}
\dit{lea}{No. My keys are on the table, \hl{next to} the lamp. But my phone isn't there.}{Non. Mes clés sont sur la table, à côté de la lampe. Mais mon téléphone n'est pas là.}
\dit{sam}{OK, I can call you\ldots{} Listen! It's \hl{under} the sofa!\note{5}}{Bon, je t'appelle… Écoute ! Il est sous le canapé !}
\dit{lea}{Under the sofa? Really?}{Sous le canapé ? Sérieusement ?}
\dit{sam}{And look: your glasses are \hl{behind} the TV!}{Et regarde : tes lunettes sont derrière la télé !}
\dit{lea}{Ha ha! My apartment is small, but I lose everything!}{Ha ha ! Mon appartement est petit, mais je perds tout !}
\end{dialogue}

\begin{comprendre}
\pointnote{4} \en{\textbf{Where's my phone?}} \en{Where's} = \en{Where is}, pour un objet. Pour
  plusieurs : \en{Where are my keys?} On répond : \en{It's…} / \en{They're…} + préposition.
\pointnote{5} \en{\textbf{I can call you}} « je peux t'appeler » : \en{can} sera étudié à
  l'unité 12. \en{Listen!} = « Écoute ! ».
\end{comprendre}

%% ---------------------------------------------------------------------
\rubrique{mots}{Les mots de l'unité}

Voici les 27 expressions de l'unité. Aux États-Unis, \en{the first floor} est le
rez-de-chaussée : on compte les étages à partir du sol.

\begin{vocabulaire}{Le logement}{8.05}
\categorie{Où j'habite}
\mot{I live in an apartment.}{aɪ ˌlɪv ɪn ən əˈpɑːrtmənt}{J'habite dans un appartement.}
\mot{a house}{ə ˈhaʊs}{une maison}
\mot{a building}{ə ˈbɪldɪŋ}{un immeuble, un bâtiment}
\mot{the first floor}{ðə ˌfɝːst ˈflɔːr}{le rez-de-chaussée (américain)}
\mot{the stairs}{ðə ˈsterz}{l'escalier}
\mot{an elevator}{ən ˈeləveɪtɚ}{un ascenseur}
\categorie{Les pièces}
\mot{a living room}{ə ˈlɪvɪŋ ˌruːm}{un salon, une salle de séjour}
\mot{a kitchen}{ə ˈkɪtʃɪn}{une cuisine}
\mot{a bedroom}{ə ˈbedruːm}{une chambre}
\mot{a bathroom}{ə ˈbæθruːm}{une salle de bains}
\mot{a balcony}{ə ˈbælkəni}{un balcon}
\mot{a garden}{ə ˈɡɑːrdən}{un jardin}
\mot{a window}{ə ˈwɪndoʊ}{une fenêtre}
\mot{a door}{ə ˈdɔːr}{une porte}
\end{vocabulaire}

\begin{vocabulaire}{Les meubles et les équipements}{8.06}
\categorie{Les meubles}
\mot{a table}{ə ˈteɪbəl}{une table}
\mot{a sofa}{ə ˈsoʊfə}{un canapé}
\mot{a bed}{ə ˈbed}{un lit}
\mot{a closet}{ə ˈklɑːzɪt}{un placard, une penderie}
\mot{a shelf}{ə ˈʃelf}{une étagère}
\mot{a mirror}{ə ˈmɪrɚ}{un miroir}
\mot{a rug}{ə ˈrʌɡ}{un tapis}
\categorie{Les équipements}
\mot{a fridge}{ə ˈfrɪdʒ}{un réfrigérateur}
\mot{a stove}{ə ˈstoʊv}{une cuisinière}
\mot{a washing machine}{ə ˈwɑːʃɪŋ məˌʃiːn}{une machine à laver}
\mot{a laundromat}{ə ˈlɑːndrəmæt}{une laverie automatique}
\categorie{Pour commenter}
\mot{on the corner}{ɑːn ðə ˈkɔːrnɚ}{au coin de la rue}
\mot{lovely}{ˈlʌvli}{charmant, ravissant}
\end{vocabulaire}

%% ---------------------------------------------------------------------
\rubrique{grammaire}{There is, there are et les prépositions}

\otIntertitre[otViolet]{1 · \en{there is} et \en{there are} : « il y a »}
\begin{center}\otTableau
\begin{tabular}{@{}l@{\hspace{7mm}}>{\enOn}l<{\enOff}@{\hspace{7mm}}>{\enOn}l<{\enOff}@{}}
 & \otTH{Singulier} & \otTH{Pluriel}\\\midrule
affirmation & There's a kitchen.             & There are two windows.\\
négation    & There isn't a garden.          & There aren't any chairs.\\
question    & Is there a bathroom?           & Are there any plants?\\
réponses    & Yes, there is. / No, there isn't. & Yes, there are. / No, there aren't.\\\bottomrule
\end{tabular}\end{center}

\begin{regle}[\en{any} au pluriel]
À la forme négative et dans les questions, on ajoute \en{any} devant un nom pluriel :
\en{There aren't any chairs} (« il n'y a pas de chaises »), \en{Are there any plants?}
(« est-ce qu'il y a des plantes ? »). Vous retrouverez \en{any} à l'unité 10.
\end{regle}

\blocaudio{8.07}{Écoutez et répétez}{Les phrases du tableau : affirmations, négations, questions
et réponses courtes.}

\otIntertitre[otViolet]{2 · Les prépositions de lieu}
\begin{tcbraster}[raster columns=4,raster equal height=rows,raster column skip=2.4mm,raster row skip=2.4mm,raster before skip=1mm,raster after skip=2mm]
\otCartePrg{blocs}{otViolet}{\en{in}}{\centering\otDessinPrepo{in}\par dans}
\otCartePrg{blocs}{otViolet}{\en{on}}{\centering\otDessinPrepo{on}\par sur}
\otCartePrg{blocs}{otViolet}{\en{under}}{\centering\otDessinPrepo{under}\par sous}
\otCartePrg{blocs}{otViolet}{\en{next to}}{\centering\otDessinPrepo{nextto}\par à côté de}
\otCartePrg{blocs}{otViolet}{\en{behind}}{\centering\otDessinPrepo{behind}\par derrière}
\otCartePrg{blocs}{otViolet}{\en{in front of}}{\centering\otDessinPrepo{infront}\par devant}
\otCartePrg{blocs}{otViolet}{\en{between}}{\centering\otDessinPrepo{between}\par entre}
\otCartePrg{blocs}{otViolet}{\en{opposite}}{\centering\vspace{2mm}\enb{\textit{The café is opposite the building.}}\par\vspace{2mm}en face de\par{\fontsize{7.6}{9}\selectfont américain : \en{across from}}}
\end{tcbraster}

\blocaudio{8.08}{Écoutez et répétez}{Sept phrases pour situer un objet : \en{The ball is in the box}…}

\begin{attention}[Les pièges des francophones]
\fauxjuste{There is two windows.}{There are two windows.}
\fauxjuste{It has a kitchen.}{There's a kitchen.}
\fauxjuste{The keys are on the bag.}{The keys are in the bag.}
\fauxjuste{next the lamp}{next to the lamp}
\par\vspace{1mm}{\fontsize{10}{12.6}\selectfont « Il y a » ne se traduit jamais par \en{it has} :
c'est toujours \en{there is} ou \en{there are}. Et à l'intérieur d'un contenant (sac, boîte,
pièce), on dit \en{in}, pas \en{on}.\par}
\end{attention}

%% ---------------------------------------------------------------------
\rubrique[Le son th de three]{sons}{Le son \phon{θ} de \en{three}}

\phon{θ} est le « frère sans voix » de \phon{ð} (unité 4) : même position de la langue, entre les
dents, mais on souffle \textbf{sans faire vibrer la voix}. On l'entend dans \en{\textbf{th}ree},
\en{\textbf{th}ink}, \en{\textbf{th}ank you}, \en{ba\textbf{th}room}, \en{mou\textbf{th}}. Les
francophones le remplacent souvent par « s », « f » ou « t » : « sank you », « free », « tree ».

\begin{paires}{8.09}{Dans chaque paire, le premier mot a le son \phon{θ}, le second un \phon{t},
\phon{s} ou \phon{f}. Écoutez, puis répétez en gardant la langue entre les dents pour le premier.}
\sonpaire{\otH{th}ree}{θriː}{trois}{\otH{t}ree}{triː}{un arbre}
\sonpaire{\otH{th}ink}{θɪŋk}{penser}{\otH{s}ink}{sɪŋk}{un évier}
\sonpaire{\otH{th}ank}{θæŋk}{remercier}{\otH{s}ank}{sæŋk}{coula}
\sonpaire{\otH{th}in}{θɪn}{mince}{\otH{s}in}{sɪn}{un péché}
\sonpaire{ma\otH{th}}{mæθ}{les maths}{ma\otH{ss}}{mæs}{la masse}
\sonpaire{pa\otH{th}}{pæθ}{un sentier}{pa\otH{ss}}{pæs}{passer}
\sonpaire{\otH{th}ought}{θɑːt}{pensa, pensé}{\otH{t}aught}{tɑːt}{enseigna, enseigné}
\sonpaire{bo\otH{th}}{boʊθ}{les deux}{boa\otH{t}}{boʊt}{un bateau}
\end{paires}

%% ---------------------------------------------------------------------
\rubrique{culture}{Se loger à New York}

\begin{paire}
\begin{culture}[\en{Apartment} ou \en{flat} ?]
Un appartement se dit \en{apartment} aux États-Unis et \en{flat} au Royaume-Uni. Un
\en{studio} est une seule pièce ; un \en{one-bedroom}, un deux-pièces (salon + chambre). Les
loyers (\en{the rent}) de New York sont parmi les plus élevés du monde.
\end{culture}
\begin{culture}[Les étages]
Aux États-Unis, \en{the first floor} est le \textbf{rez-de-chaussée}, et \en{the second floor} le
premier étage. Au Royaume-Uni, comme en France, on parle de \en{the ground floor}, puis du
\en{first floor}.
\end{culture}
\end{paire}

\begin{culture}[La vie d'immeuble]
Les vieux immeubles de Brooklyn, en grès brun, s'appellent des \en{brownstones}. Beaucoup n'ont ni
ascenseur ni machine à laver : on lave son linge à la laverie du quartier, le \en{laundromat}. Dans
les grands immeubles, un gardien (\en{the super}, pour \en{superintendent}) s'occupe des réparations.
\end{culture}

%% ---------------------------------------------------------------------
\rubrique{pratique}{Exercices}

\exercice[ecrire,niveau=1]{Dans quelle pièce ?}{Complétez avec le nom d'une pièce.}
\begin{questions}
\item You sleep in the \trou[34mm]{bedroom}.
\item You cook in the \trou[34mm]{kitchen}.
\item You take a shower in the \trou[34mm]{bathroom}.
\item You watch TV in the \trou[34mm]{living room}.
\item Flowers and plants are in the \trou[34mm]{garden}.
\end{questions}

\exercice[ecrire,niveau=1]{\en{there is} ou \en{there are} ?}{Complétez avec \en{there's},
\en{there are}, \en{there isn't} ou \en{there aren't}, d'après la description de l'appartement de Léa.}
\begin{questions}
\item In Léa's apartment, \trou[30mm]{there's} a small living room.
\item \trou[30mm]{There are} two big windows.
\item \trou[30mm]{There isn't} a washing machine.
\item \trou[30mm]{There's} a laundromat on the corner.
\item \trou[30mm]{There are} plants on the fire escape.
\item \trou[30mm]{There aren't} any chairs on the fire escape.
\end{questions}

\exercice[ecrire,niveau=2]{Questions et réponses}{Répondez par une réponse courte, d'après le
dialogue 1.}
\begin{questions}
\item Is there a bathroom? \hfill \trou[50mm]{Yes, there is.}
\item Is there a washing machine? \hfill \trou[50mm]{No, there isn't.}
\item Are there two windows? \hfill \trou[50mm]{Yes, there are.}
\item Is there a laundromat? \hfill \trou[50mm]{Yes, there is.}
\item Are there plants? \hfill \trou[50mm]{Yes, there are.}
\item Is there a big kitchen? \hfill \trou[50mm]{No, there isn't.}
\end{questions}

\exercice[lire,ecrire,niveau=1]{Où est la balle ?}{Regardez chaque dessin et complétez avec la
bonne préposition.}
\begin{tcbraster}[raster columns=3,raster equal height=rows,raster column skip=3mm,raster row skip=3mm,raster before skip=1mm]
\otCartePrepo{on}{The ball is}{on}{the box.}
\otCartePrepo{under}{The ball is}{under}{the table.}
\otCartePrepo{in}{The ball is}{in}{the box.}
\otCartePrepo{behind}{The ball is}{behind}{the box.}
\otCartePrepo{between}{The ball is}{between}{the boxes.}
\otCartePrepo{nextto}{The ball is}{next to}{the box.}
\end{tcbraster}

\exercice[lire,niveau=2]{L'annonce}{Lisez cette annonce, puis indiquez si chaque phrase est vraie (V)
ou fausse (F).}
\begin{anglais}
\begin{tcolorbox}[otBase,colback=white,colframe=otRule,boxrule=0.5pt,fontupper=\fontsize{10.8}{14.4}\selectfont]
{\otSemi\fontsize{11}{13}\selectfont BROOKLYN — Sunny one-bedroom apartment}\par\vspace{1.4mm}
There's a living room with two windows, a kitchen, a bedroom and a small bathroom. There isn't a
washing machine, but there's a laundromat on the corner. The apartment is on the third floor.
There isn't an elevator. Rent: \$2,400 a month.
\par\vspace{1mm}{\fontsize{8.6}{10}\selectfont\color{otMuted}\foreignlanguage{french}{\textit{sunny} : ensoleillé · \textit{rent} : loyer · \textit{a month} : par mois}}
\end{tcolorbox}
\begin{questions}[2]
\item There are two bedrooms. \hfill\trou[9mm]{F}
\item There's a washing machine. \hfill\trou[9mm]{F}
\item The laundromat is on the corner. \hfill\trou[9mm]{V}
\item There's an elevator. \hfill\trou[9mm]{F}
\item There are two windows in the living room. \hfill\trou[9mm]{V}
\item The apartment is on the third floor. \hfill\trou[9mm]{V}
\end{questions}
\end{anglais}

\exercice[ecrire,niveau=2]{Corrigez l'erreur}{Chaque phrase contient une erreur. Récrivez-la correctement.}
\begin{questions}
\item There is three bedrooms. \hfill \trou[70mm]{There are three bedrooms.}
\item It has a garden behind the house. \hfill \trou[70mm]{There's a garden behind the house.}
\item My keys are next the lamp. \hfill \trou[70mm]{My keys are next to the lamp.}
\item Is there any plants? \hfill \trou[70mm]{Are there any plants?}
\item The phone is on my bag. \hfill \trou[70mm]{The phone is in my bag.}
\end{questions}

\exercice[ecouter,niveau=1,piste=8.10]{\phon{θ} ou pas ?}{Écoutez. Pour chaque numéro, la voix dit
un seul des deux mots, deux fois. Cochez le mot entendu.}
\begin{questions}[4]
\item \choix{three}\quad\choix{tree}
\item \choix{think}\quad\choix{sink}
\item \choix{thank}\quad\choix{sank}
\item \choix{thin}\quad\choix{sin}
\item \choix{math}\quad\choix{mass}
\item \choix{path}\quad\choix{pass}
\item \choix{thought}\quad\choix{taught}
\item \choix{both}\quad\choix{boat}
\end{questions}

\exercice[ecouter,ecrire,niveau=2,piste=8.11]{Les affaires de Jack}{Jack ne trouve rien ce matin.
Emma l'aide. Écoutez et notez où sont ses affaires.}
\begin{center}\otTableau
\begin{tabular}{@{}l@{\hspace{14mm}}l@{}}
\otTH{Objet} & \otTH{Où ?}\\\midrule
{\otSemi 1 · \enb{his keys}}     & \enb{\trou[60mm]{on his desk, next to the computer}}\\[1mm]
{\otSemi 2 · \enb{his umbrella}} & \enb{\trou[60mm]{behind the door}}\\[1mm]
{\otSemi 3 · \enb{his phone}}    & \enb{\trou[60mm]{under his papers}}\\[1mm]
{\otSemi 4 · \enb{his coffee}}   & \enb{\trou[60mm]{in his hand}}\\\bottomrule
\end{tabular}\end{center}

\exercice[ecouter,ecrire,niveau=2,piste=8.12]{Dictée à trous}{Mrs.\ Lopez présente son immeuble.
Écoutez et complétez.}
\begin{anglais}
\begin{tcolorbox}[otBase,colback=white,colframe=otRule,boxrule=0.5pt,fontupper=\fontsize{11}{19}\selectfont]
Our building is old. There \trou{are} six apartments. \trou{There's} no elevator, but there
\trou{are} stairs! Your apartment is \trou{on} the third floor. My apartment is on the first floor,
\trou{next} to the door. There's a small garden \trou{behind} the building. And there \trou{is}
a coffee shop \trou{opposite} the building!
\end{tcolorbox}
\end{anglais}

\exercice[ecrire,niveau=2]{Traduisez}{Traduisez en anglais.}
\begin{questions}[1][french]
\item Il y a une cuisine et deux chambres.\lignes{1}
\item Il n'y a pas d'ascenseur.\lignes{1}
\item Est-ce qu'il y a une salle de bains ? — Oui. \textit{(réponse courte)}\lignes{1}
\item Où sont mes clés ? — Elles sont sur la table.\lignes{1}
\item Le chat est sous le lit.\lignes{1}
\item La laverie est au coin de la rue.\lignes{1}
\item J'habite dans un appartement au deuxième étage.\lignes{1}
\item Mon téléphone est dans mon sac, à côté de mes lunettes.\lignes{1}
\end{questions}

\exercice[ecrire,niveau=3]{Chez vous}{Décrivez votre logement en six phrases : le type de logement,
les pièces, ce qu'il y a et ce qu'il n'y a pas, et votre pièce préférée.}
\lignes{6}

\exercice[ecouter,parler,niveau=3,piste=8.13]{À vous de parler !}{Emma vous pose des questions sur
votre logement. Répondez-lui \textbf{à voix haute} pendant chaque pause.}
\begin{center}\otTableau
\begin{tabular}{@{}>{\otSemi\fontsize{8.6}{10}\selectfont\color{otNavy}}l@{\hspace{3mm}}>{\enOn}p{70mm}<{\enOff}@{\hspace{6mm}}p{78mm}@{}}
 & \otTH{Emma dit\ldots} & \otTH{Vous\ldots}\\\midrule
1. & Do you live in a house or an apartment? & répondez.\\
2. & Is there a garden? & répondez par une réponse courte.\\
3. & How many bedrooms are there? & dites combien il y a de chambres.\\
4. & Where's your kitchen? & situez votre cuisine (à côté de…).\\
5. & It sounds lovely! & remerciez-la.\\\bottomrule
\end{tabular}\end{center}

%% ---------------------------------------------------------------------
\rubrique{bilan}{Faites le point}

\begin{essentiel}
\begin{minipage}[t]{0.52\linewidth}
\otEssTitre{Les phrases clés}
\renewcommand{\arraystretch}{1.32}
\begin{tabular}{@{}>{\enOn\otSemi\fontsize{9.4}{11}\selectfont}l<{\enOff}@{\hspace{3mm}}>{\fontsize{9.2}{11}\selectfont\color{white!78!otNavy}}l@{}}
I live in an apartment.      & J'habite un appartement.\\
There's a kitchen.           & Il y a une cuisine.\\
There are two windows.       & Il y a deux fenêtres.\\
Is there a bathroom?         & Il y a une salle de bains ?\\
Where's my phone?            & Où est mon téléphone ?\\
It's under the sofa.         & Il est sous le canapé.\\
\end{tabular}
\end{minipage}\hfill
\begin{minipage}[t]{0.42\linewidth}
\otEssTitre{Il y a}
{\enb{\fontsize{9.4}{12}\selectfont there's + singulier · there are + pluriel}\par
{\fontsize{9}{11.4}\selectfont\color{white!78!otNavy}négation, question : \en{any} au pluriel}\par}
\otEssTitre{Prépositions}
{\enb{\fontsize{9.2}{11.6}\selectfont in · on · under · next to · behind · in front of · between · opposite}\par}
\otEssTitre{Prononciation}
{\fontsize{9.2}{11.6}\selectfont \phon{θ} : langue entre les dents, sans la voix : \en{three}, \en{bathroom}\par}
\end{minipage}
\end{essentiel}

\begin{jesais}
\item décrire mon logement et nommer les pièces et les meubles ;
\item dire ce qu'il y a avec \en{there's} et \en{there are} ;
\item poser la question \en{Is there…?} / \en{Are there any…?} et y répondre ;
\item situer un objet avec huit prépositions de lieu ;
\item demander où se trouve un objet : \en{Where's…?} \en{Where are…?} ;
\item prononcer \phon{θ} sans le confondre avec \phon{s}, \phon{f} ou \phon{t}.
\end{jesais}

\clearpage
\begin{testexpress}
\qcm[b]{« Il y a deux chambres » se dit :}{\en{There is two bedrooms.}}{\en{There are two bedrooms.}}{\en{It has two bedrooms.}}
\qcm[a]{« Il n'y a pas d'ascenseur » se dit :}{\en{There isn't an elevator.}}{\en{There aren't elevator.}}{\en{It isn't an elevator.}}
\qcm[c]{\en{Is there a garden?} — }{\en{Yes, it is.}}{\en{Yes, there's.}}{\en{Yes, there is.}}
\qcm[b]{« sous le lit » se dit :}{\en{on the bed}}{\en{under the bed}}{\en{in the bed}}
\qcm[a]{« à côté de la lampe » se dit :}{\en{next to the lamp}}{\en{next the lamp}}{\en{behind the lamp}}
\qcm[c]{Les clés sont dans le sac :}{\en{The keys are on the bag.}}{\en{The keys are at the bag.}}{\en{The keys are in the bag.}}
\qcm[b]{\en{the first floor}, aux États-Unis, c'est :}{le premier étage}{le rez-de-chaussée}{le grenier}
\qcm[a]{« un appartement » se dit au Royaume-Uni :}{\en{a flat}}{\en{a studio}}{\en{an apartment}}
\qcm[c]{Le son \phon{θ} de \en{three} se prononce :}{comme un « s »}{comme un « t »}{langue entre les dents, sans la voix}
\qcm[b]{« Où sont mes lunettes ? » se dit :}{\en{Where's my glasses?}}{\en{Where are my glasses?}}{\en{Where my glasses are?}}
\end{testexpress}

\noindent{\fontsize{9.6}{12}\selectfont\color{otMuted}Niveau visé (CECRL, niveau A1) : « peut
décrire simplement son logement et situer des objets ». Les corrigés de l'unité sont
page~\pageref{cor-u8-1} ; les textes des pistes d'exercices, page~\pageref{trans-8.10}.\par}
